\documentclass[10pt,a4paper]{amsart}
\usepackage{amsmath,amssymb,amsthm,graphicx}
% For commutative diagrams we use Xy-pic
\usepackage[all]{xy}

% We use 2cell for 2-commutative diagrams.
\xyoption{2cell}
% Theorem environments.
%
\theoremstyle{plain}
\newtheorem{theorem}[subsection]{Theorem}
\newtheorem{proposition}[subsection]{Proposition}
\newtheorem{lemma}[subsection]{Lemma}

\theoremstyle{definition}
\newtheorem{definition}[subsection]{Definition}
\newtheorem{example}[subsection]{Example}
\newtheorem{exercise}[subsection]{Exercise}
\newtheorem{situation}[subsection]{Situation}

\theoremstyle{remark}
\newtheorem{remark}[subsection]{Remark}
\newtheorem{remarks}[subsection]{Remarks}

\numberwithin{equation}{subsection}

% Macros
%
\def\lim{\mathop{\mathrm{lim}}\nolimits}
\def\colim{\mathop{\mathrm{colim}}\nolimits}
\def\Spec{\mathop{\mathrm{Spec}}}
\def\Hom{\mathop{\mathrm{Hom}}\nolimits}
\def\Ext{\mathop{\mathrm{Ext}}\nolimits}
\def\SheafHom{\mathop{\mathcal{H}\!\mathit{om}}\nolimits}
\def\SheafExt{\mathop{\mathcal{E}\!\mathit{xt}}\nolimits}
\def\Sch{\mathit{Sch}}
\def\Mor{\mathop{\mathrm{Mor}}\nolimits}
\def\Ob{\mathop{\mathrm{Ob}}\nolimits}
\def\Sh{\mathop{\mathit{Sh}}\nolimits}
\def\NL{\mathop{N\!L}\nolimits}
\def\CH{\mathop{\mathrm{CH}}\nolimits}
\def\proetale{{pro\text{-}\acute{e}tale}}
\def\etale{{\acute{e}tale}}
\def\QCoh{\mathit{QCoh}}
\def\Ker{\mathop{\mathrm{Ker}}}
\def\Im{\mathop{\mathrm{Im}}}
\def\Coker{\mathop{\mathrm{Coker}}}
\def\Coim{\mathop{\mathrm{Coim}}}

% Boxtimes
%
\DeclareMathSymbol{\boxtimes}{\mathbin}{AMSa}{"02}

%
% Macros for moduli stacks/spaces
%
\def\QCohstack{\mathcal{QC}\!\mathit{oh}}
\def\Cohstack{\mathcal{C}\!\mathit{oh}}
\def\Spacesstack{\mathcal{S}\!\mathit{paces}}
\def\Quotfunctor{\mathrm{Quot}}
\def\Hilbfunctor{\mathrm{Hilb}}
\def\Curvesstack{\mathcal{C}\!\mathit{urves}}
\def\Polarizedstack{\mathcal{P}\!\mathit{olarized}}
\def\Complexesstack{\mathcal{C}\!\mathit{omplexes}}
% \Pic is the operator that assigns to X its picard group, usage \Pic(X)
% \Picardstack_{X/B} denotes the Picard stack of X over B
% \Picardfunctor_{X/B} denotes the Picard functor of X over B
\def\Pic{\mathop{\mathrm{Pic}}\nolimits}
\def\Picardstack{\mathcal{P}\!\mathit{ic}}
\def\Picardfunctor{\mathrm{Pic}}
\def\Deformationcategory{\mathcal{D}\!\mathit{ef}}
\usepackage{cleveref}
\crefname{originallemma}{original lemma}{original lemmas}
\Crefname{originallemma}{Original lemma}{Original lemmas}
\makeatletter
\newlabel{more-algebra-lemma-perfect}{{0658}{source}{}{sourcepointer-1}{}}
\newlabel{more-algebra-lemma-perfect@cref}{{[originallemma][0][]0658}{[1][0][]source}{}{}{}}
\makeatother
\crefname{originallemma}{original lemma}{original lemmas}
\Crefname{originallemma}{Original lemma}{Original lemmas}
\makeatletter
\newlabel{sites-cohomology-lemma-locally-constant}{{094G}{source}{}{sourcepointer-2}{}}
\newlabel{sites-cohomology-lemma-locally-constant@cref}{{[originallemma][0][]094G}{[1][0][]source}{}{}{}}
\makeatother
\title{On a connected scheme a bounded below constructible complex with locally constant cohomology is constant on an etale covering over finite projective modules}
\author{The Stacks Project Authors}
\date{stacks-project:etale-cohomology.tex:local-snapshot-20261011; GNU FDL 1.2 or later}
\begin{document}
\maketitle
\noindent\textit{Editorial note.} All delivered mathematical text is the author's own, copied from the pinned Stacks Project chapter file etale-cohomology.tex; no statement or proof step was rewritten, reordered or removed.\par\medskip
\noindent\textit{Editorial note.} The counted proof is the author's complete argument: the proof is written out in full: the trivialising covering, the grouping into disjoint open pieces by isomorphism class of the local model with the reason for each step, the connectedness conclusion, and the perfectness criterion making the model finite projective.\par\medskip
\noindent\textit{Editorial note.} The label lemma-connected-ctf-locally-constant names the counted statement, and the proof environment immediately following it in the source is delivered verbatim as its proof.\par\medskip
\noindent\textit{Editorial note.} The counted statement is carried by the author's own lemma environment in the chapter source.\par\medskip
\noindent\textit{Editorial note.} The statement and proof are copied from the pinned chapter file; the two cited inputs are kept as external pointers with their canonical Stacks tags.\par\medskip
\noindent\textit{Editorial note.} The preamble is the chapter's own native preamble excerpt already used for this source: the Xy-pic setup (native preamble lines 14-18) and the native theorem declarations with the mathematical macros (native preamble lines 154-217).\par\medskip
\noindent\textit{Allowed standard prerequisites (editorial).}\begin{itemize}
\item The constructible finite type derived category with locally constant cohomology. Same chapter, the definition of the full subcategory of complexes with constructible cohomology sheaves and the lemma identifying the constants for a locally constant sheaf of Lambda-modules. Supplies the hypothesis on K and the local constancy description used at the start of the proof.
\item Perfectness of the local models. More on Algebra, the lemma characterising perfect complexes over Noetherian rings through the finite type and boundedness properties. Supplies that the module complexes M\_i\textasciicircum{}bullet are perfect, hence replaceable by finite complexes of finite projective modules.
\end{itemize}\medskip
\ifdefined\hypertarget\hypertarget{sourcepointer-1}{}\fi
\noindent\textit{Source reference.} Original lemma 0658: sources/stacks-chow/tags.txt, source lines 7892--7892. Prerequisite pointer invoked at proof line(s) 13875; it names a standard fact of a sibling chapter that the counted argument uses, and every step of the counted proof built on it is retained verbatim in the delivered body.\par\medskip
\ifdefined\hypertarget\hypertarget{sourcepointer-2}{}\fi
\noindent\textit{Source reference.} Original lemma 094G: sources/stacks-chow/tags.txt, source lines 11549--11549. Prerequisite pointer invoked at proof line(s) 13863; it names a standard fact of a sibling chapter that the counted argument uses, and every step of the counted proof built on it is retained verbatim in the delivered body.\par\medskip
\begin{lemma}
\label{lemma-connected-ctf-locally-constant}
Let $X$ be a connected scheme. Let $\Lambda$ be a Noetherian ring.
Let $K \in D_{ctf}(X_\etale, \Lambda)$ have locally constant cohomology sheaves.
Then there exists a finite complex of finite projective $\Lambda$-modules
$M^\bullet$ and an \'etale covering $\{U_i \to X\}$ such that
$K|_{U_i} \cong \underline{M^\bullet}|_{U_i}$ in $D(U_{i, \etale}, \Lambda)$.
\end{lemma}

\begin{proof}
Choose an \'etale covering $\{U_i \to X\}$ such that $K|_{U_i}$
is constant, say $K|_{U_i} \cong \underline{M_i^\bullet}_{U_i}$
for some finite complex of finite $\Lambda$-modules $M_i^\bullet$.
See Cohomology on Sites, Lemma
\ref{sites-cohomology-lemma-locally-constant}.
Observe that $U_i \times_X U_j$ is empty if $M_i^\bullet$
is not isomorphic to $M_j^\bullet$ in $D(\Lambda)$.
For each complex of $\Lambda$-modules $M^\bullet$ let
$I_{M^\bullet} =
\{i \in I \mid M_i^\bullet \cong M^\bullet\text{ in }D(\Lambda)\}$.
As \'etale morphisms are open we see that
$U_{M^\bullet} = \bigcup_{i \in I_{M^\bullet}} \Im(U_i \to X)$
is an open subset of $X$. Then $X = \coprod U_{M^\bullet}$ is a disjoint
open covering of $X$. As $X$ is connected only one $U_{M^\bullet}$
is nonempty. As $K$ is in $D_{ctf}(X_\etale, \Lambda)$ we see that $M^\bullet$
is a perfect complex of $\Lambda$-modules, see
More on Algebra, Lemma \ref{more-algebra-lemma-perfect}.
Hence we may assume $M^\bullet$ is a finite complex of finite projective
$\Lambda$-modules.
\end{proof}

\end{document}
